% Standalone source; no external bibliography or figure files are required.
% Compile twice with pdflatex to resolve the contents and cross-references.
\documentclass[11pt,a4paper]{article}
\usepackage[T1]{fontenc}
\usepackage[utf8]{inputenc}
\usepackage[margin=25mm]{geometry}
\usepackage{amsmath,amssymb}
\usepackage{booktabs,array,longtable,tabularx}
\usepackage{enumitem}
\usepackage{microtype}
\usepackage{url}
\usepackage[hidelinks]{hyperref}
\urlstyle{tt}
\setlength{\emergencystretch}{2em}
\setlength{\parskip}{0.35em}
\setlength{\parindent}{0pt}
\setlist{itemsep=0.2em,topsep=0.35em}
\newcolumntype{Y}{>{\raggedright\arraybackslash}X}
\newcolumntype{P}[1]{>{\raggedright\arraybackslash}p{#1}}
\newcommand{\code}[1]{\texttt{\detokenize{#1}}}
\newcommand{\stopgrad}{\operatorname{sg}}
\newcommand{\R}{\mathbb{R}}
\DeclareMathOperator*{\argmin}{arg\,min}
\DeclareMathOperator*{\argmax}{arg\,max}

\title{Methodology of the EAGLE Replica\\
\large Implementation, Experimental Protocol, and Replication Fidelity}
\author{Source-based technical description}
\date{18 November 2026}

\begin{document}
\maketitle

\begin{abstract}
This document describes the implementation of Efficient Autoscaling via
Generalizable multi-agent reinforcement Learning Environment (EAGLE) in the
local EAGLE-replica repository. It follows the terminology of Monaco et al.\ and
distinguishes the published design from the behavior implemented in the source.
The methodology covers trace preparation, network partitioning and controller
placement, Mininet/Ryu operation, graph embedding, average-reward collaborative
Actor--Critic learning, global reward estimation, pruning-based consensus,
checkpointing, and the available evaluation workflows. A separate section
identifies absent or partial components and assesses whether their faithful
replication is possible from the supplied code and paper. No new training,
emulation, or performance results are claimed.
\end{abstract}

\tableofcontents
\newpage

\section{Methodology}
\label{sec:methodology}

\subsection{Scope, evidence, and terminology}
\label{sec:scope}

The primary scientific reference is \emph{Adaptive SDN Autoscaling via
Generalizable Multi-Agent Reinforcement Learning With EAGLE}, published in
\emph{IEEE Transactions on Network and Service Management}, volume~23,
pages~6507--6522, 2026~\cite{paper}. The implementation described here is the local
working tree inspected on 18 November 2026, whose Git HEAD was
\code{37699c8af206e1a2daf3faae37e085efbaa9e5f7}. Source files and settings determine
the implemented procedure; explanatory guides and the notebook provide
workflow context. Generated datasets, networks, models, and evaluation outputs
are treated as local artifacts rather than as specifications of the original
experiments. Most of these output directories are excluded from version
control~\cite{code-snapshot}.

Throughout this document, a \emph{community}, \emph{subnet}, and \emph{domain}
refer to a switch partition managed by a logical controller. An \emph{agent}
means the learning entity associated with a logical controller slot, such as
\code{c0}. A \emph{physical controller} is a running Ryu process. This distinction
matters because a surviving Ryu process can serve a failed controller's logical
slot without changing the number or order of learning agents. The implementation
calls a controllable switch a \emph{support switch}; by default, one such switch
is assigned to each domain. References to numbered equations or figures
\emph{in the paper} refer to the original publication, whereas unqualified
equation references refer to this document.

The paper supplies the objective and the principal algorithms. The code supplies
additional operational choices, including trace-window sampling, time
stretching, normalization references, optimizer selection, training budgets,
and confidence-interval construction. These choices are reported explicitly
instead of being attributed to the paper. The reconstruction assessment in
Section~\ref{sec:missing} uses the supplied local source and paper; the upstream
repository mentioned in the paper was not independently inspected.

\subsection{Network model, autoscaling objective, and control loop}

Following the paper's network model, let $G=(V,E,w)$ be an undirected switch
graph, where $V$ denotes switches, $E$ denotes inter-switch links, and $w_l$
denotes a link's capacity. Hosts attach to switches and generate directed
traffic demands. The network is divided into $N$ communities $c_i$, each
associated with an agent. The binary variable $\Psi_u$ denotes whether switch
$u$ is available for forwarding. The paper seeks a throughput--power trade-off,
written in its Equation~(4) as
\begin{equation}
  \max_{\Psi}\; Th_{\mathrm{norm}}-\alpha PC_{\mathrm{norm}}.
  \label{eq:objective}
\end{equation}
The paper's capacity, flow-conservation, switch/link-availability, and
single-path constraints provide the optimization motivation. The replica
implements an RL control loop and a shortest-path routing module, rather than
solving that mathematical program at each decision~\cite{paper,code-routing}.

For the implemented action space, only the support switches are controlled by
RL. Let $u_i$ be domain $i$'s support switch and $a_t^i\in\{0,1\}$ its decision
at step $t$. Then $\Psi_{u_i}=a_t^i$; other reachable switches remain available.
With three domains, the joint action has $2^3$ possible configurations. The
joint policy factorizes into local actors,
\begin{equation}
  \pi_\theta(a_t\mid s_t)=\prod_{i=1}^{N}
  \pi_{\theta^i}(a_t^i\mid s_t),
\end{equation}
where every actor receives the same concatenated, normalized state. The
long-term learning objective follows the globally averaged, average-reward
formulation in the paper's Equation~(19):
\begin{equation}
 J(\theta)=\lim_{T\rightarrow\infty}\frac{1}{T}
 \mathbb{E}_{\pi_\theta}\left[
   \sum_{t=0}^{T-1}\frac{1}{N}\sum_{i=1}^{N}r_{t+1}^i
 \right].
 \label{eq:average-objective}
\end{equation}
Finite-episode cumulative rewards are reporting quantities derived from this
setting, rather than estimates obtained by explicitly evaluating the limit.

The closed loop has the following stages: controllers collect telemetry;
active resources form weighted virtual subgraphs; each subgraph is embedded;
the embeddings are concatenated in controller order and normalized; the actors
select binary actions; Ryu updates forwarding; the environment waits for
acknowledgments and post-action telemetry; and rewards and next states are
constructed. During training, batches additionally trigger local gradient
updates and a consensus step. During evaluation, the neural networks remain
in inference mode~\cite{code-env,code-learning}.

\subsection{Software decomposition and three-runtime execution}

The implementation separates privileged emulation, the legacy controller
dependencies, and numerical learning into three runtimes
(Table~\ref{tab:runtimes}). Mininet is launched with system Python and root
privileges; Ryu uses its Python~3.7 virtual environment; preparation, training,
and evaluation use the model's Python~3.13 environment. The controller
requirements pin Ryu~4.34 and NetworkX~2.6.3, while the model uses native
PyTorch and newer numerical libraries~\cite{code-runtime}.

\begin{table}[htbp]
\centering\small
\caption{Runtime responsibilities and primary entry points.}
\label{tab:runtimes}
\begin{tabularx}{\textwidth}{P{0.17\textwidth}YY}
\toprule
Runtime & Responsibilities & Primary source \\
\midrule
System Python / Mininet & OVS bridges, host namespaces, shaped links,
traffic replay, delivery measurements &
\path{mininet_nets/launch_topology.py}, \path{mininet_nets/traffic.py} \\
Python 3.7 / Ryu & OpenFlow ownership, route installation, peer presence,
active topology, counter polling &
\path{ryu_controller/launch_controllers.py}, \path{ryu_controller/controller.py} \\
Python 3.13 / model & Prepared inputs, embeddings, learners, consensus,
checkpoints, frozen-policy evaluations &
\path{data_preparation/}, \path{network_preparation/}, \path{model/} \\
\bottomrule
\end{tabularx}
\end{table}

Ryu must start before Mininet, with the identical network JSON supplied to both
launchers. For example, the documented runtime commands are
\begin{quote}\small
\path{ryu_controller/.venv/bin/python}
\path{ryu_controller/launch_controllers.py}
\code{--scenario networks/v0.json}

\code{sudo python3} \path{mininet_nets/launch_topology.py}
\code{--scenario networks/v0.json}.
\end{quote}
Model entry points instead accept a filename such as \code{--network v0.json},
resolved beneath \path{networks/}. These are execution instructions for the
methodology, not operations performed while preparing this document.

The Ryu supervisor creates a Unix control socket and writes its address to
the root \path{.socket}; Mininet writes its address to
\path{mininet_nets/.socket}. The numerical runtime discovers these addresses,
with explicit overrides available. Mininet's socket handles status, reset,
traffic, and finish operations. The Ryu supervisor handles status and restoration
of its selected processes. Changing topology requires manually restarting the
launchers; the notebook preserves that operational boundary~\cite{code-env,code-runtime}.

\subsection{Traffic-trace preparation}
\label{sec:traffic}

\subsubsection{Source captures and complete one-second windows}

The default inputs are the Digital Corpora capture
\path{5gb-tcp-connection.pcap.gz} and the Megalodon Content Server capture
\path{megalodon.pcapng}. The parser supports Ethernet PCAP and PCAPNG,
including compressed input. It reads timestamps and original frame lengths;
the optional analysis workflow reports per-direction traffic statistics and
one-second throughput histories~\cite{code-data}.

Preparation groups packets into one-second windows according to their original
Unix timestamps. If $t_{\min}$ and $t_{\max}$ are the first and last packet
times, retained window starts are the integers in
$[\lceil t_{\min}\rceil,\lfloor t_{\max}\rfloor)$. Thus partial boundary
windows are excluded, while empty complete windows are retained. Digital
Corpora windows are shuffled with a seeded generator and split at
$\lfloor4W/5\rfloor$, producing approximately 80\% training and 20\% test
windows. The split unit is a time window, rather than an entire TCP connection.

Megalodon is used for test traffic. Its windows are classified into an
\emph{isolated} directed IP pair,
\code{192.168.1.1 -> 172.28.190.238}, and the \emph{remaining} traffic. Both
phase-specific window sets enter the pool. The subsequent allocation is
shuffled, so the generated traffic does not preserve the source capture's
complete chronology or its original end-to-end sessions.

\subsubsection{Host-pair allocation, minimum horizon, and time stretching}

For $H$ hosts, preparation enumerates $P=H(H-1)$ ordered source--destination
pairs. In the default allocation, hosts have conventional identifiers
\code{h1}, \code{h2}, and so forth, unique locally administered MAC addresses,
and IPv4 addresses in \code{10.0.0.0/8}. Host metadata is written into the
dataset manifests and subsequently consumed by network preparation.

Let $W$ be the available windows in a split, $q$ the number assigned to each
directed pair, $\sigma$ the time-stretch factor, and $\tau$ the configured
decision interval. Preparation first considers $q_0=\lfloor W/P\rfloor$.
Only if
\begin{equation}
  \left\lfloor\frac{q_0\sigma}{\tau}\right\rfloor<n_{\min}
\end{equation}
does it increase the allocation to approximately
$q=\lceil n_{\min}\tau/\sigma\rceil$, with an additional floating-point
rounding check. Otherwise it retains $q_0$, including any horizon longer than
the minimum. The defaults are $n_{\min}=20$ and $\tau=5$~s. The factor
\code{Config.TRAFFIC_SCALE_FACTOR} defaults to $H$ when unset, and must be at
least one. The resulting dataset duration is $D=q\sigma$~\cite{code-data}.

Windows are drawn without replacement within a pool traversal, and a fresh
shuffled draw is used when further windows are required. This permits reuse
when the available scaled duration cannot meet the requested minimum, while
keeping Digital Corpora's training and test pools separate. The pool's seed
also incorporates the split and source name. With nine hosts, for example,
twelve source windows per direction become 108~s, supplying 21 complete
five-second decisions. The minimum is checked \emph{after} scaling.

\subsubsection{Prepared PCAPs and synthetic TCP workloads}

Each directed pair has two possible representations derived from the same
allocated windows. In the PCAP representation, packets retain their
within-window ordering and have timestamps
\begin{equation}
  t' = \sigma\bigl(j+(t\bmod 1\;\mathrm{s})\bigr),
\end{equation}
where $j$ is the allocated window index. Ethernet and supported network-layer
addresses are rewritten for the assigned hosts. The implementation updates
IPv4 header checksums and relevant transport checksums, including handling
supported fragmentation and IPv6 extension-header cases. Packet and frame-byte
counts are retained, and each pair's manifest records the first packet's offset
and its prepared capture path~\cite{code-data}.

In the synthetic representation, each nonempty pair/window becomes a managed
TCP transfer. If the source window contains $B_j$ original frame bytes, its
transfer size is the rounded positive byte count and its target offered rate
is approximately
\begin{equation}
  b_j=\frac{8B_j}{\sigma}\quad\mathrm{bit/s}.
\end{equation}
The job starts at $j\sigma$, occupies a nominal window of $\sigma$ seconds,
and has timeout $\max(60,\lceil\sigma\rceil+5)$ seconds. Original frame-byte
volume is therefore used to set \emph{application payload} volume; this is a
synthetic workload construction, not a byte-for-byte reproduction of the
original transport sessions. Time stretching preserves assigned volume while
reducing offered rate. Empty windows produce no TCP job.

Preparation writes separate train/test manifests, supporting \code{--skip-pcaps}
for synthetic-only output and \code{--skip-train} for test-only output. It stages
new datasets in a temporary directory and rejects an existing destination
instead of silently overwriting it. Manifests record durations, host metadata,
scaling, decision counts, and seeds where produced by the current generator.

\subsection{Traffic conditions and curriculum domains}
\label{sec:curriculum-data}

The code implements the paper's source/target traffic terminology through
named manifests. Table~\ref{tab:traffic} distinguishes capture format from
load level. In particular, T2 is the Megalodon synthetic dataset
\code{mg_t0}, rather than a file named \code{dc_t2}~\cite{code-data,code-datasets}.

\begin{table}[htbp]
\centering\small
\caption{Prepared source and target traffic conditions. Multipliers apply to
already time-stretched synthetic traffic.}
\label{tab:traffic}
\begin{tabularx}{\textwidth}{P{0.21\textwidth}P{0.24\textwidth}Y}
\toprule
Paper terminology & Manifest stem & Implemented construction \\
\midrule
Source T0 / TrEnv 1 & \code{dc_t0} & Digital Corpora synthetic jobs \\
Source capture & \code{dc_pcap} & Remapped Digital Corpora PCAPs \\
Target T1 & \code{dc_t1} & $0.5$ times Digital Corpora T0 volume/rate \\
Target T2 & \code{mg_t0} & Megalodon synthetic jobs \\
Alternative capture & \code{mg_pcap} & Remapped Megalodon PCAPs \\
Target T3 & \code{dc_t3} & $1.5$ times Digital Corpora T0 volume/rate \\
Target T4 & \code{dc_t4} & $2$ times Digital Corpora T0 volume/rate \\
TrEnvs 2--4 & \code{dc_t0_trenv2} through \code{dc_t0_trenv4} & $1.5$ times
traffic originating in the first, second, or third controller domain \\
TrEnv 5 & \code{dc_t0_trenv5} & $1.5$ times traffic originating at all hosts \\
\bottomrule
\end{tabularx}
\end{table}

Congestion variants require a partitioned network with exactly three distinct
controllers. Host ownership is obtained from attachment switches; controller
list order determines the three regional environments. The generator verifies
that dataset directions equal the network's complete set of ordered host pairs.
Only originating traffic is multiplied for a regional hotspot, including both
intra-domain and inter-domain destinations. Window durations and the global
time-stretch factor remain unchanged.

These environments operationalize the paper's \emph{Domain Incremental
Learning}: the topology and action semantics stay fixed while the traffic
distribution changes. Their names specify intended congestion conditions;
the generator does not verify that a particular link actually exceeds
capacity. For the same split, TrEnv~5 and T3 implement the same global
$1.5$ multiplier, so they are not independent traffic constructions.

\subsection{Topology construction, balanced Louvain partitioning, and placement}
\label{sec:topologies}

\subsubsection{Random and Topology Zoo switch graphs}

The random generator creates a connected simple graph with exactly the
requested switch and link counts. It first constructs a randomized spanning
tree and then samples additional absent edges. This procedure fixes graph
size, but does not constrain the degree sequence to equal V0's. Topology Zoo
processing removes marked hosts and self-loops, collapses directed/parallel
links to simple undirected links, retains the largest connected component,
and deterministically renumbers its switches. Original Zoo identifiers are
retained as provenance~\cite{code-network}.

For newly generated random/Zoo scenarios, inter-switch capacities are sampled
as integer Mbps uniformly from 58 to 100 using the selected seed. Host-link
capacities use the same range. Capacities are stored in bit/s and also serve
as the weights for partitioning. Source V0 capacities are preserved from its
source JSON instead of being resampled.

\subsubsection{Balanced Louvain heuristic}

The partitioner implements the two phases of Louvain: local node moves and
contraction of communities into supernodes. Its modularity calculations include
a resolution parameter $\xi$. Writing $k_u=\sum_v w_{uv}$ and
$m=\sum_{\{u,v\}\in E}w_{uv}$, the corresponding modularity is
\begin{equation}
 Q_\xi=\frac{1}{2m}\sum_{u,v}
 \left(w_{uv}-\xi\frac{k_uk_v}{2m}\right)
 \mathbf{1}\{c(u)=c(v)\}.
\end{equation}
This is the resolution-adjusted form of the paper's Equation~(11). Nodes are
visited in a seeded shuffled order. Candidate destinations are neighboring
communities, and a move must exceed a small positive gain threshold.

The move-capacity parameter is $\gamma=\lceil |V|/N\rceil$ by default, or an
explicit \code{LOUVAIN_CAPACITY}. Supernodes retain their original members, so
capacity is measured in original switches at every contraction level. A move
whose destination would exceed $\gamma$ is rejected. If the proposed partition's
largest and smallest nonempty communities differ by more than $\epsilon$, the
gain is penalized by $\rho$. Defaults are $\epsilon=2$, $\rho=0.01$, and
$\xi=0.1$~\cite{code-network,code-settings}.

The algorithm retains a Louvain level with at least the requested number of
domains, splits any disconnected communities into connected components, and
uses greedy modularity merging on the quotient graph when necessary to obtain
exactly $N$ domains. The move-capacity condition applies to the balanced
Louvain phase; it is not reimposed during this final merge. Thus final domains
need not meet the earlier capacity or balance condition. Both pre-merge and
final partitions are recorded and plotted.

\subsubsection{Controller locations, support switches, and host attachment}

For each final community $c_i$, placement chooses a member minimizing its mean
unweighted hop distance to community switches:
\begin{equation}
 p_i=\argmin_{p\in c_i}\frac{1}{|c_i|}
       \sum_{u\in c_i}d_G(p,u).
 \label{eq:placement}
\end{equation}
Distances are computed on the full switch graph, and deterministic node
ordering breaks ties. Random/Zoo controller slots are ordered by their
placement switches. This implements the minimum-hop objective associated with
the paper's \emph{controller placement problem}, using unit switch loads.
The location is scenario metadata; the default runtime endpoints are local
loopback addresses~\cite{code-network,code-routing}.

Support-switch selection ranks non-articulation switches by descending degree,
then identifier, and uses backtracking with connectivity checks. The selected
set must permit switch-disable combinations without disconnecting the
remaining switch network; selected switches also retain a neighbor when the
other selected switches are disabled. Host attachment avoids selected support
switches, favors low-degree peripheral switches, and fills available slots
round-robin across domains. Each eligible access switch receives at most one
host. Generation rejects requests exceeding the available access slots.

The scenario generator can record several \code{controllable_switches} per
domain, but the runtime and actor use only \code{support_switch}, which is the
first selected entry. The default \code{SWITCHES_PER_DOMAIN=1} matches that
implemented binary-action interface.

\subsubsection{Preserved source V0 and expanded V3}

V0 is reconstructed from \path{network_preparation/source_v0.json}. This file
preserves 13 switches, 16 inter-switch links, source port assignments,
capacities, community ownership, and nine host-attachment positions. The
generator replaces host identities and addresses with dataset metadata, then
recomputes placement locations within the preserved domains. Logical slots
remain ordered as
\begin{equation}
 (\text{\code{c0}},\text{\code{c1}},\text{\code{c2}})
 \longleftrightarrow
 (\text{\code{s12}},\text{\code{s11}},\text{\code{s13}})
\end{equation}
for support-switch actions. The resulting placement locations in the inspected
V0 artifact are \code{s6}, \code{s11}, and \code{s7}; placement and support
switch are distinct concepts.

V3 preserves V0 and adds three switches and six inter-switch links:
\code{s14} belongs to \code{c0} and connects to \code{s1}/\code{s6};
\code{s15} belongs to \code{c2} and connects to \code{s3}/\code{s7}; and
\code{s16} belongs to \code{c1} and connects to \code{s5}/\code{s9}.
New capacities use the configured seed and the 58--100~Mbps range.
Existing host attachments and support-switch order are retained~\cite{code-network}.

\begin{table}[htbp]
\centering\small
\caption{Inspected local V0--V4 artifacts. Links count simple undirected
inter-switch edges, excluding host links. These are artifact observations,
not a claim of identity with all original paper topologies.}
\label{tab:topologies}
\begin{tabularx}{\textwidth}{@{}lrrrY@{}}
\toprule
Network & Switches & Links & Hosts & Construction / qualification \\
\midrule
V0 & 13 & 16 & 9 & Preserved source graph and domains \\
V1 & 13 & 16 & 9 & Random connected graph; degree sequence differs from V0 \\
V2 & 13 & 18 & 9 & Random connected graph with more links \\
V3 & 16 & 22 & 9 & Explicit extension of V0 \\
V4 & 149 & 193 & 50 & GtsCe; the paper's Table~III reports 386 links \\
\bottomrule
\end{tabularx}
\end{table}

The bundled GtsCe GraphML itself contains 149 nodes and 193 undirected edges.
The factor-of-two difference from the paper's 386 may reflect a directed-arc
count, but that interpretation is not established by the supplied material.
The inspected V4's 50-host configuration is a local choice; the generator
accepts other compatible host-count datasets. Its final domain sizes are
2, 94, and 53, illustrating that post-Louvain merging can override balance.
V0's preserved domain sizes are 5, 3, and 5, rather than a newly computed
Louvain partition.

\subsection{Scenario identity and input compatibility}

Version-1 network JSON describes controllers, endpoint addresses, switch
datapath identifiers, host addresses, link endpoints, explicit port numbers,
capacities, telemetry intervals, and provenance. A SHA-256 digest of the raw
JSON bytes identifies the scenario across the model, Ryu, and Mininet.
Consequently, even formatting changes alter the fingerprint. A different
traffic dataset can be selected without changing network identity
\cite{code-env,code-datasets,code-routing}.

Dataset resolution uses the selected network's host count, the requested
train/test split, and the manifest stem. Host identifiers, MAC addresses, and
IPv4 addresses must match the network. Jobs must have distinct identifiers,
valid distinct endpoints, and start times inside the declared duration.
Dataset fingerprints include manifest bytes and referenced prepared-PCAP
digests. The traffic runner verifies PCAP digests again before replay. For
older manifests without explicit hosts, the loader can recover addresses from
prepared IPv4 packets or conventional host identifiers; current generation
writes the metadata directly.

Controller order is part of the checkpoint contract. The saved encoder's
controller identifier tuple must exactly match the scenario, including order,
and all agent states must be present. This allows reuse across different switch
and host counts with the same logical slots, but does not support arbitrary
changes in agent number or slot order.

\subsection{Mininet data plane and the reaction phase}
\label{sec:forwarding}

Mininet builds OVS kernel switches in secure fail mode, host namespaces with
configured addresses, and \code{TCLink} links shaped to the scenario capacities.
Switches enable OpenFlow~1.3 and~1.5 and connect to the configured controllers.
Static ARP entries support host communication. The Ryu application selects
OpenFlow~1.5 by default and contains serialization paths for both versions
\cite{code-runtime,code-routing}.

The \emph{configured graph} contains all scenario switches and hosts. The
\emph{active graph} contains switches reachable through their elected physical
owners, except support switches whose action is zero. It retains only links
whose relevant switch ports are up and removes isolated hosts. An action
therefore removes or restores a switch and its incident links in the forwarding
representation. Ryu reconciles forwarding rules with this graph. The action
path does not stop an OVS process, lower a physical interface, or measure a
hardware power transition.

For each ordered host pair, the routing planner enumerates hop-shortest
end-to-end paths in the active graph. Among equal-cost paths, it prefers those
through the source domain's support switch, if available. A stable CRC32C hash
selects one of the remaining candidates. Rules match source and destination
Ethernet MAC addresses, and each pair follows one selected path. Multiple TCP
connections between that host pair therefore share the same forwarding choice.
This is a deterministic host-pair realization of equal-cost path selection;
the implementation does not install OpenFlow ECMP groups~\cite{code-routing}.

Only changed EAGLE-owned rules are added or deleted, identified by a dedicated
cookie. A new master rebuilds the relevant entries. OpenFlow barriers track
completion of rule dispatch, and telemetry is reset after topology or ownership
changes. This ties the reaction phase to subsequent observations of the new
forwarding configuration.

\subsection{Controller collaboration, ownership, and telemetry}
\label{sec:telemetry}

\subsubsection{Peer presence and physical controller failover}

Each Ryu process exposes its presence, connected switches and active ports,
resource actions, and revision information through HTTP. Controllers poll peers
at a default interval of 0.3~s and expire stale presence after 2~s. Once all
configured peers and switches have been seen during startup, a live primary
serves its own logical slot. A failed primary's slot is assigned to the first
live peer in scenario order. OpenFlow master/slave role requests implement
the resulting ownership~\cite{code-routing}.

Support-switch decisions propagate as revisioned state. Revisions combine a
nanosecond timestamp and controller identifier, and newer revisions supersede
older ones. The HTTP client first contacts a slot's own endpoint and then tries
other endpoints, allowing a surviving controller to serve that slot's statistics
or action requests. Logical slots remain fixed. Clock synchronization is an
operational assumption for multi-host timestamps and revision ordering.

\subsubsection{Coherent counter windows and virtual subgraphs}

OpenFlow port counters and host-attachment flow counters are polled with a
default interval of 1~s. Multipart replies are correlated by datapath and
transaction identifier. A usable rate sample requires two complete rounds from
the same topology generation. For an observed counter $C_t$, the increment is
$C_t-C_{t-1}$ when nonnegative; a lower counter is treated as a reset and the
new value becomes the increment. Rates use the measured separation between
poll-round start times, not an assumed one-second duration.

For a directional port increment $\Delta B$ over $\Delta t$, the utilization
weight is
\begin{equation}
 u_l=\begin{cases}
 0, & 8\Delta B/(\Delta t\,w_l)<0.01,\\
 \min\{1,8\Delta B/(\Delta t\,w_l)\}, & \text{otherwise}.
 \end{cases}
\end{equation}
The undirected observation edge takes the maximum of the available directional
weights. Thus values below one percent are suppressed and saturation is clipped
to one~\cite{code-routing}.

A slot's \emph{virtual subgraph} contains its active local switches and every
active edge incident on those switches, along with the edges' endpoints.
Consequently, host nodes and neighboring switches from other domains can appear
in the graph, and adjacent slots' observation graphs can overlap. This is more
specific than merely taking an induced graph on the domain's switches.

Delivered traffic is attributed to the destination host's attachment switch
using the EAGLE forwarding-flow counters, keyed by directed host MAC pair.
These counts are an OpenFlow byte-based delivery proxy, rather than a direct
application-payload goodput measurement. Local switch/port counts and active
switch egress rates provide the power-model inputs. Inactive forwarding
resources are excluded from the modeled usage even though their emulated
datapaths can remain connected.

\subsection{Joint decisions, sample freshness, and episode boundaries}
\label{sec:episodes}

\code{JointEnvironment} exposes one binary action per slot through a
\code{MultiDiscrete} space and a dictionary containing \code{global_obs} as its
observation. Rewards are vectors in scenario order. Episode reset checks the
Mininet fingerprint, requests restoration of supervised Ryu processes when a
supervisor socket is available, finishes earlier traffic, and applies an all-on
joint action. It waits for post-action telemetry before scheduling the next
dataset. A configurable warm-up is then followed by a fresh observation
\cite{code-env}.

At each decision, actions are submitted concurrently and successful
acknowledgments are retained while unavailable slots are retried. An
acknowledgment carries the accepted action, scenario, revision, physical owner,
and acceptance time. The environment waits for at least the nominal
\code{decision_seconds} from the start of the step and obtains new snapshots.
A snapshot set must satisfy all of the following: each sample window begins
after the latest action acceptance; all views contain the acknowledged actions
and revisions; snapshot times differ by at most three poll intervals; and
the oldest sample is no more than three poll intervals old. The default
sample timeout is 15~s. Failure to obtain a valid joint action or observation
raises an error rather than manufacturing a transition.

The default horizon is
\begin{equation}
 K=\left\lfloor\frac{D-t_{\mathrm{warmup}}}{\tau}\right\rfloor,
 \label{eq:horizon}
\end{equation}
with at least one required decision. Reset synchronization, delayed telemetry,
and learning updates can extend actual elapsed time, so Equation~\eqref{eq:horizon}
specifies a nominal decision count rather than a strict wall-clock stopping
time. Step durations are measured. Each transition reports
\code{terminated=False}; reaching $K$ sets \code{truncated=True}. Learning
therefore bootstraps across the time limit instead of treating it as an
absorbing network state.

After a successful trial, a default 5~s drain period permits outstanding
deliveries under the final forwarding configuration. Rewards and modeled power
are accumulated over decision steps; deliveries completing during the drain
can contribute to FCT. Reset, warm-up, and drain durations do not enter the
reported decision-time mean power~\cite{code-trials}.

\subsection{Implemented reward and switch power model}
\label{sec:reward}

For agent $i$, let $\mathcal{F}_{i,t}$ be the directed host-pair entries present
in its delivered-byte sample, and let $B_{f,t}$ denote a byte increment over
sample duration $\Delta t$. The implemented throughput statistic is
\begin{equation}
 Th_{i,t}=\begin{cases}
 \displaystyle\frac{1}{|\mathcal{F}_{i,t}|}
      \sum_{f\in\mathcal{F}_{i,t}}\frac{8B_{f,t}}{\Delta t},
      & |\mathcal{F}_{i,t}|>0,\\
 0, & \text{otherwise}.
 \end{cases}
 \label{eq:throughput}
\end{equation}
This averages available destination-domain host-pair counters, rather than
all individual transport connections. The globally averaged reward gives equal
weight to agents; it need not equal a global per-flow average when domains
contain different numbers of sampled pairs~\cite{code-reward}.

Let $n^{\mathrm{sw}}_{i,t}$ and $n^{\mathrm{port}}_{i,t}$ be the active local
switch and port counts, and $b^{\mathrm{egress}}_{u,t}$ a switch's active-port
transmit rate in bit/s. The modeled local power is
\begin{equation}
 PC_{i,t}=30n^{\mathrm{sw}}_{i,t}+n^{\mathrm{port}}_{i,t}
       +0.05\sum_{u\in c_i}\frac{b^{\mathrm{egress}}_{u,t}}{10^6}
       \quad\mathrm{W}.
 \label{eq:power}
\end{equation}
An active inter-switch link contributes a port at each active switch endpoint;
a host-facing link contributes its switch-side port. This implements the
paper's static device/port component and rate-dependent component in
Equation~(3), with numerical coefficients chosen by the local settings.

With $Th_{\mathrm{ref}}=8\times10^6$ bit/s,
$PC_{\mathrm{ref}}=200$ W, and $\alpha=0.5$, the local reward is
\begin{equation}
 r_{t+1}^{i}=
 \min\left\{1,\frac{Th_{i,t}}{Th_{\mathrm{ref}}}\right\}
 -\alpha\min\left\{1,\frac{PC_{i,t}}{PC_{\mathrm{ref}}}\right\}.
 \label{eq:reward}
\end{equation}
References are fixed configuration values, not recomputed per topology or
relative to an all-on baseline. The reward lies between $-0.5$ and $1$ for
nonnegative counters under these defaults. Both terms can saturate, particularly
on large domains. Total reported network power sums the local modeled powers.
An episode's cumulative reward is
\begin{equation}
 C_{\mathrm{episode}}=\sum_{t=0}^{K-1}
                  \frac{1}{N}\sum_{i=1}^{N}r_{t+1}^i.
 \label{eq:cumulative}
\end{equation}
It is a raw sum, not divided by episode length or min--max normalized.

\subsection{FeatherGraph state representation and source normalization}
\label{sec:embedding}

\subsubsection{Fixed-size graph fingerprints}

The replica implements a graph-level FEATHER/FeatherGraph-style descriptor
directly using NetworkX, NumPy, and SciPy. The paper's path length $g$ maps to
\code{embedding.order}; its number of evaluation points $e$ maps to
\code{embedding.eval_points}. The implementation uses characteristic-function
features and sparse neighborhood propagation, rather than sampling individual
walks at runtime~\cite{paper,code-embedding}.

For an observation graph, nodes are sorted and implicit unit self-loops are
added. Let $A$ be its weighted adjacency, including a unit diagonal, and let
$d_u$ be the degree used by the implementation. It selects weighted degree
only if NetworkX considers every edge weighted. Because the inserted self-loops
have implicit unit weights without explicit weight attributes, ordinary
telemetry graphs use \emph{unweighted} degree: a self-loop counts twice in
$d_u$ but once in $A_{uu}$. This convention is explicitly covered by a numerical
unit fixture. Therefore the propagation matrix below should not be silently
replaced by a conventional row-normalized weighted adjacency.

With unweighted clustering coefficient $\kappa_u$, the node descriptors are
$x_u=(\log(1+d_u),\kappa_u)$. Evaluation points are
\begin{equation}
 (\omega_1,\ldots,\omega_e)=\operatorname{linspace}(0.01,2.5,e).
\end{equation}
For $e=1$, the sole point is $0.01$. Define the $4e$-dimensional row
\begin{equation}
 \Phi_u=\bigl[\{\cos(\omega_jx_{u,k})\}_{k,j},
                   \{\sin(\omega_jx_{u,k})\}_{k,j}\bigr],
 \qquad k\in\{1,2\},\quad j\in\{1,\ldots,e\}.
\end{equation}
Let $M=\operatorname{diag}(d_u)^{-1}A$. Starting from $F^{(0)}=\Phi$,
the graph fingerprint concatenates mean-pooled propagated features:
\begin{equation}
 F^{(q)}=MF^{(q-1)},\qquad
 z_i=\biggl[\frac{1}{|V_i'|}\sum_{u\in V_i'}F^{(q)}_u
                                      \biggr]_{q=1}^{g}.
 \label{eq:feather}
\end{equation}
An empty observation graph yields a zero vector. Each local embedding has
$4ge$ components; the concatenated vector
\begin{equation}
 Z_t=[z_{1,t};\ldots;z_{N,t}]\in\R^{4geN}
\end{equation}
uses scenario controller order. Defaults $g=e=1$ and $N=3$ give four features
per slot and twelve global features. Mean pooling removes dependence on node
ordering, while link utilization and active-graph structure affect the
descriptor. The encoder is unsupervised and is not a trainable GNN.

\subsubsection{Calibration on source training conditions}

Before learning, the default \code{source} normalization mode runs one all-on
calibration episode for each selected training condition. It collects raw
post-step global embeddings and obtains componentwise minimum $m_j$ and
maximum $M_j$, requiring at least two finite calibration vectors. The normalized
input is
\begin{equation}
 s_{t,j}=\begin{cases}
 100\,(Z_{t,j}-m_j)/(M_j-m_j), & M_j>m_j,\\
 0, & M_j=m_j.
 \end{cases}
 \label{eq:normalization}
\end{equation}
Values are not clipped to $[0,100]$: target-domain or action-induced values can
lie outside the calibration range. Nonfinite outputs are replaced by zero,
and observations are converted to float32. An optional \code{identity} mode
uses sanitized raw float32 embeddings instead~\cite{code-embedding}.

Calibration bounds and slot order are saved in both training checkpoints and
inference artifacts. Evaluation reuses these source bounds; it does not fit
normalization to the target test traffic. Curriculum training calibrates on
all selected source TrEnvs. The fixed feature size permits topology-size
transfer only while the logical agent interface remains compatible; proximity
in embedding space is an empirical diagnostic rather than a guarantee of
zero-shot performance.

\subsection{Actor--Critic agents and global reward estimation}
\label{sec:learning}

\subsubsection{Neural-network architecture and action selection}

Each agent maintains an Actor, Critic, and Reward estimator, parameterized by
$\theta^i$, $v^i$, and $\lambda^i$, respectively, following the paper's
collaborative AC-MARL terminology. All are multilayer perceptrons with hidden
widths 256, 128, and 64 and ReLU activations. They use the same architectural
pattern but have independent parameter tensors. The actor maps the global state
to two logits, the critic maps it to one scalar state value, and the estimator
maps the state concatenated with the joint binary action to one scalar reward
estimate~\cite{code-learning}.

With the default interface, actor and critic inputs have twelve components;
the estimator input has fifteen. Actors initialize linear weights uniformly
in $[-0.04,0.04]$ and biases to zero. Critic and estimator weights use Xavier
uniform initialization, with biases left at PyTorch's linear-layer defaults.
All three networks use Adam, with learning rates
$\beta_\theta=10^{-4}$, $\beta_v=10^{-3}$, and
$\beta_\lambda=10^{-3}$. These rate values and hidden widths match the paper's
training setup; the explicit Adam implementation is a local optimizer choice.

Training samples $a_t^i$ from the categorical distribution defined by actor
logits. Each slot has a separately seeded CPU action generator, including when
the networks run on CUDA. Inference uses the greedy logit argmax for each slot,
with no policy updates. Model seeds initialize network parameters by slot, and
the action-generator state is checkpointed. The \code{auto} device setting
selects CUDA if available; explicit CPU and CUDA selections are supported.

\subsubsection{Reward-estimator pretraining}

The paper defines $\bar R(s,a)$ as the mean of local rewards and proposes a
feed-forward estimator trained through local targets and parameter consensus
(paper Equations~(20)--(22)). In the replica, estimator learning takes place in
a distinct pretraining phase. For batch $\mathcal B$, agent $i$ minimizes
\begin{equation}
 L_\lambda^i=\frac{1}{|\mathcal B|}\sum_{t\in\mathcal B}
 \left[\bar R(s_t,a_t;\lambda^i)-r_{t+1}^i\right]^2.
 \label{eq:estimator-loss}
\end{equation}
The action is concatenated as $N$ float-valued binary entries, rather than as
a joint-action one-hot vector. Actors continue sampling their initialized
policies during this phase, but the actor and critic are not updated.
Estimator parameter consensus follows each batch~\cite{code-learning,code-consensus}.

After pretraining, estimator inference supplies the actor's global-reward proxy
and estimator weights remain fixed throughout policy learning. This follows
the pretrain-then-inference setup stated in the paper's Section~VI-A. Its
Section~V also presents a general update with estimator sharing in the critic
step; the local implementation resolves that presentation by separating the
two phases. Parameter averaging of nonlinear estimators should be understood
as the implemented approximation to global reward estimation, not as an exact
arithmetic average of every peer's predicted reward.

\subsubsection{Average-reward critic and actor updates}

During policy learning, each collected transition contains the global state,
joint action, reward vector, next global state, and termination/truncation
flags. Each learner reads its own component of the reward vector. The
coordinator nevertheless possesses the full vector because it constructs and
dispatches the joint batch~\cite{code-learning,code-workers}.

Let $h_t=1-\mathbf{1}\{\mathrm{terminated}_t\}$ and $\mu^i$ be agent $i$'s
running average-reward estimate. Using pre-update critic values, define
\begin{align}
 y_t^i &= r_{t+1}^i-\mu^i+h_tV(s_{t+1};v^i),\\
 \delta_t^i &= y_t^i-V(s_t;v^i),\\
 \widetilde\delta_t^i &= \bar R(s_t,a_t;\lambda^i)-\mu^i
             +h_tV(s_{t+1};v^i)-V(s_t;v^i).
 \label{eq:td}
\end{align}
The last two expressions correspond to the local state-value TD-error and
globally averaged TD-error in the paper's Equations~(25) and~(29), with the
implementation's termination mask made explicit. There is no discounted-return
factor: the next-state coefficient is one in this average-reward formulation.

The implemented losses are
\begin{align}
 L_v^i &= \frac{1}{|\mathcal B|}\sum_{t\in\mathcal B}
       \left[V(s_t;v^i)-\stopgrad(y_t^i)\right]^2,\\
 L_\theta^i &= -\frac{1}{|\mathcal B|}\sum_{t\in\mathcal B}
       \log\pi_{\theta^i}(a_t^i\mid s_t)
       \stopgrad(\widetilde\delta_t^i).
 \label{eq:ac-losses}
\end{align}
Here $\stopgrad$ denotes a detached target. The actor advantage is computed
before the critic optimizer step and is held fixed during the actor update.
Adam minimizes these losses; consequently, the code is a minibatch
semi-gradient realization of the paper's update direction, rather than a
literal scalar-step implementation of its Equations~(26) and~(28).

The running average is initialized to zero and updated once per batch:
\begin{equation}
 \widetilde\mu^i=(1-\beta_v)\mu^i+
       \beta_v\frac{1}{|\mathcal B|}\sum_{t\in\mathcal B}r_{t+1}^i.
 \label{eq:mu}
\end{equation}
The critic parameters and $\widetilde\mu^i$ then enter consensus. Estimator
pretraining does not update $\mu^i$. Since normal episode horizons truncate
rather than terminate the MDP, the final transition still bootstraps.

\subsection{Consensus-based collaborative learning and pruning}
\label{sec:consensus}

\subsubsection{Synchronous mixing and worker execution}

The default learner pool spawns one Python process per logical controller;
\code{--workers local} keeps the same learners in the coordinator process.
The coordinator distributes observations and batches, requests local updates,
collects contributions, and sends each learner the other learners'
contributions. Local embedding can also be dispatched to these workers.
Communication uses multiprocessing pipes and trusted pickle messages, rather
than a peer-to-peer transport between Ryu hosts~\cite{code-workers}.

All fixed slots participate in every learning round with uniform weight
$1/N$. Actors remain local. Pretraining shares only the estimator;
policy learning shares only the critic and the scalar average reward, because
the estimator is frozen. For a shared tensor $W$, the effective update at
receiver $i$ is
\begin{equation}
 W^i_{r+1}=\frac{1}{N}\widetilde W^i_r+
          \frac{1}{N}\sum_{j\ne i}P_{p_r}(\widetilde W^j_r),
 \label{eq:mix}
\end{equation}
where $P_{p_r}$ is the outgoing pruning operator. Scalars are averaged without
pruning. The round number, component names, parameter keys, and tensor shapes
are checked before mixing. A learner process failure or timeout aborts the
round. Physical Ryu failover does not remove a learner slot from this pool.

\subsubsection{Local unstructured pruning of outgoing copies}

Following the paper's pruning-based strategy, each agent copies its shared
weights and applies per-layer magnitude pruning before transmission. The
first linear layer and all biases remain dense. For every eligible later
weight tensor with $K_W$ entries, the smallest
$\operatorname{round}(p_rK_W)$ absolute values are removed, using stable
ordering for ties. Surviving float32 values and a packed bitmask are serialized
with shape metadata. The receiver reconstructs missing entries as zero
\cite{code-consensus}.

The sender's local tensor is not sparsified when constructing a message.
Subsequent consensus still changes local parameters through
Equation~\eqref{eq:mix}; preserving the local outgoing source does not imply
that the mixed model is numerically unchanged. Peer weights are not renormalized
coordinatewise when a contribution has been pruned.

Three modes are implemented: no pruning ($p_r=0$), static pruning ($p_r=0.7$
from the first round), and dynamic pruning. With dynamic warm-up $r_0=5$ and
ramp length $R=20$, the exact schedule is
\begin{equation}
 p_r=\begin{cases}
 0, & r\leq r_0,\\
 0.3+0.4\min\{1,(r-r_0-1)/R\}, & r>r_0.
 \end{cases}
 \label{eq:pruning}
\end{equation}
Thus round~6 uses 30\% and round~26 reaches 70\%. Communication rounds count
batch updates continuously across pretraining and policy learning, rather
than restarting at each stage or corresponding one-to-one to episodes.

The learning log records an intended all-to-all parameter traffic count
\begin{equation}
 B_r^{\mathrm{peer}}=(N-1)\sum_{i=1}^{N}|\mathrm{packet}_{i,r}|.
\end{equation}
This includes serialized contribution payloads but is not a measurement of
network bandwidth, physical packet overhead, or all exchanged state/action
information. A separate pool counter tracks serialized coordinator/worker
traffic in process mode.

\subsection{Training schedule, defaults, and reproducibility artifacts}
\label{sec:training}

Three training modes select source manifests automatically: \code{synthetic}
uses Digital Corpora \code{dc_t0}; \code{pcaps} uses \code{dc_pcap}; and
\code{curriculum} uses T0 followed by TrEnvs~2--5. Every mode follows the same
outer procedure: source calibration when enabled, estimator pretraining,
then Actor--Critic learning with frozen estimators~\cite{code-training}.

Pretraining cycles through the selected profiles over its episode budget.
Policy-learning episodes are divided evenly among curriculum stages, with
remainders assigned to earlier stages. The default five pretraining episodes
and thirty policy episodes therefore provide one estimator episode and six
policy episodes per curriculum profile. Calibration runs one additional
all-on episode per profile. Synthetic-only and PCAP-only modes each have one
profile. Curriculum requires at least five policy episodes.

Transitions are accumulated in non-overlapping batches of eight. A full batch
triggers one local-update/consensus round, and a final partial batch is flushed
at each episode boundary. The dataset schedule is restarted for each episode;
different action samples and runtime timing can change observations, but a new
source trace is not resampled on every episode. Training uses controller
telemetry and disables per-flow delivery capture to reduce measurement work.

\begin{table}[htbp]
\centering\small
\caption{Principal defaults from \texttt{config.py} and
\texttt{settings/default.json}. Fixed choices not specified by the paper are
implementation settings rather than recovered original hyperparameters.}
\label{tab:defaults}
\begin{tabularx}{\textwidth}{YY}
\toprule
Parameter & Default \\
\midrule
Seed / domains / RL-controlled switches & 42 / 3 / one support switch per domain \\
Decision / warm-up / drain & 5~s / 0~s / 5~s \\
Telemetry sample timeout & 15~s \\
FeatherGraph $g$, $e$ / global width & 1, 1 / 12 at three slots \\
Calibration / estimator / policy episodes & 1 per profile / 5 total / 30 total \\
Actor, critic, estimator rates & $10^{-4}$, $10^{-3}$, $10^{-3}$ \\
Optimizer / transition batch & Adam / 8 \\
Reward $\alpha$ / throughput / power references & 0.5 / 8~Mbps / 200~W \\
Switch / port / dynamic power & 30~W / 1~W / 0.05~W per Mbps \\
Dynamic pruning & 5 dense rounds, then 30\% to 70\% over 20 round increments \\
Comparison episodes / repetitions & 5 / 10 \\
\bottomrule
\end{tabularx}
\end{table}

Settings files merge into the defaults, with unknown keys rejected and
numeric constraints checked. Named overrides select no or static pruning.
Command-line options can override episode budgets, seed, worker mode, device,
and control sockets. A typical source model is trained through
\path{model/.venv/bin/python} \code{-m model.train --network v0.json}
\code{--traffic synthetic --output dc_t0}. The usual comparison names are
\code{dc_t0}, \code{dc_pcaps}, and \code{dc_full} for the synthetic, PCAP, and
curriculum roles, respectively.

A per-name advisory file lock prevents concurrent training writers. Atomic
episode-boundary checkpoints are saved under \path{models/checkpoints/};
named inference artifacts are saved under \path{models/}. A training checkpoint
contains all three networks per slot, optimizer state, average reward,
action-generator state, encoder bounds, cumulative communication round,
completed episode count, stage schedule, reward history, settings, seed, and
scenario/dataset fingerprints. Inference artifacts retain actors, critics,
encoder, and settings; the estimator is not needed for greedy deployment.

Resume checks that network bytes, datasets, settings, seed, and training mode
match. It restarts at an episode boundary, rather than preserving a partially
executed traffic replay or minibatch. Model loading maps tensors through CPU
and supports a later device selection. Reproducible numerical state does not
make emulator timing or real TCP execution bitwise deterministic.

The policy reward CSV has columns \code{episode,cumulative_reward}; calibration
and pretraining are logged separately from that curve. A PNG accompanies the
completed training curve. Run manifests record the full scenario, hashes,
settings, seed, encoder, Python version, and major numerical-library versions.
JSON-lines events retain reset observations, transitions, telemetry,
acknowledgments, learning diagnostics, and episode summaries
\cite{code-checkpoints,code-trials}.

\subsection{Traffic execution and flow-completion measurements}
\label{sec:fct}

\subsubsection{Managed TCP transfers}

Synthetic jobs are executed by one asynchronous receiver process per
participating destination host and one asynchronous scheduled-sender process
per participating source host. Receivers listen on TCP port~20000. Once
receivers and optional captures are ready, the runner schedules a common
start approximately one second in the future. Each job waits for its absolute
scheduled start, opens a connection, sends identifying metadata, and writes
payload in chunks of up to 65,536 bytes with rate-based pacing
\cite{code-runtime}.

For a completed synthetic job $f$, the completion measure is
\begin{equation}
 FCT_f=t^{\mathrm{receiver\ complete}}_f
             -t^{\mathrm{sender\ attempt}}_f.
 \label{eq:managed-fct}
\end{equation}
The origin is the actual sender attempt, rather than the requested schedule or
the receiver's first payload byte. This includes connection establishment and
the transfer's pacing. Receiver completion requires receipt of the prescribed
payload size. Missing or incomplete delivery yields an undefined FCT.

\subsubsection{Prepared-PCAP replay and payload coverage}

PCAP jobs invoke \code{tcpreplay --multiplier=1} on the source host interface.
Their start includes the manifest's first-packet offset, preserving initial
idle time in the prepared schedule. The packet timestamps themselves already
include time stretching; replay does not apply that factor a second time.
This is open-loop frame replay, whereas synthetic TCP jobs create live
transport connections.

For comparison evaluations, \code{tcpdump} records source and destination
interfaces, filtering by the job's source/destination MAC addresses and TCP.
The analysis groups packets by directional IP-address/port four-tuple. It
builds expected unique TCP payload sequence ranges from the prepared replay
file, accounting for overlaps and retransmissions. Observed received ranges
are intersected with those expected ranges. A replay flow completes only
when all expected unique payload bytes for its four-tuple have been observed
at the destination. Its start time is the first corresponding source-capture
packet, and its finish is the packet that completes expected payload coverage
\cite{code-trials}.

This measures delivery of the \emph{prepared replay's} payload, not necessarily
completion of a full original application connection. TCP flows with no
payload supply no FCT sample, and non-TCP traffic does not enter this FCT
analysis. After draining, only records with defined FCTs are included in
means and written as completed-flow events. Incomplete deliveries are omitted
instead of receiving a timeout penalty. Consequently, FCT must be interpreted
as conditional on successful completion. Per-trial packet captures are
temporary runtime files and are removed when the next traffic directory is
prepared; the implementation does not archive every capture permanently.

\subsection{Frozen-policy evaluation, baselines, and reported metrics}
\label{sec:evaluation}

\subsubsection{Matched experiments and aggregation}

The common comparison runner loads three distinct named inference models and
an all-on baseline. The baseline holds every support switch on throughout
the episode. The three default model roles compare training regimes of the
same EAGLE replica, rather than different published autoscaling algorithms.
Each uses its own saved source encoder and greedy actions. The test dataset is
identical across methods, and each trial resets forwarding to all-on before
running its selected policy~\cite{code-evaluation}.

For repetition $b$, all methods use seed $s+b$, starting at $b=0$. Defaults
are ten repetitions and five episodes per method per condition. Repetitions
reuse the prepared traffic schedule; they do not independently resample ten
different source captures. Methods run sequentially, with the baseline first.
Within a repetition, episode FCT means are weighted by completed-flow count;
episode powers are weighted by measured decision time.

At the episode level, with step durations $\Delta t_t$, total modeled power is
reported as
\begin{equation}
 \overline{PC}=
 \frac{\sum_t\Delta t_t\sum_i PC_{i,t}}{\sum_t\Delta t_t}.
 \label{eq:mean-power}
\end{equation}
This holds each sampled power estimate over the measured decision duration;
it is not an integral from a physical power meter. The episode FCT mean is
the arithmetic mean over completed delivery records, including those completed
during the drain period.

\subsubsection{FCT gain, power saving, and efficiency score}

For model $m$ and matched repetition $b$, the runner computes the paper's
FCT gain and power saving (paper Equations~(32)--(33)):
\begin{align}
 G_{m,b}&=1-\frac{FCT_{m,b}}{FCT_{\mathrm{baseline},b}},\\
 PS_{m,b}&=1-\frac{\overline{PC}_{m,b}}
                   {\overline{PC}_{\mathrm{baseline},b}}.
 \label{eq:gains}
\end{align}
FCT gain remains a fraction; power-saving bars display $100PS$. For each
dataset/topology condition, let
$G_{\max}=\max_m\operatorname{mean}_bG_{m,b}$ over the three selected models.
The implemented relative efficiency score is
\begin{equation}
 E_{m,b}=\frac{G_{m,b}/G_{\max}}{1-PS_{m,b}}.
 \label{eq:efficiency}
\end{equation}
This realizes the paper's Equation~(34), but the normalization set is the
replica's three training variants. It is not normalized against the paper's
five competing frameworks. Scores can exceed one and are meaningful only
with the corresponding FCT gain and power saving~\cite{code-plots}.

No completed baseline flows, nonpositive baseline denominators, nonpositive
$G_{\max}$, or a power-saving denominator of zero produce undefined applicable
metrics. Undefined values are retained as empty numeric results rather than
converted to favorable scores. Negative FCT gains are allowed when defined.

Error bars are 90\% percentile bootstrap intervals. The code takes 5,000
resamples of the valid repetition-level metric values, computes each resampled
mean, and reports the 5th/95th percentiles. Pairing occurs when each model is
compared to its same-seed baseline before bootstrapping. $G_{\max}$ is held
fixed during this bootstrap rather than re-estimated within each resample.
Thus the statistical unit is a repetition, not an individual flow or an
individual episode.

\subsection{The six implemented evaluation workflows}
\label{sec:eval-workflows}

All evaluation outputs are rooted at \path{eval/}. Topology changes require
manual service restarts, Ryu first. The workflows below describe supported
entry points; their presence does not imply that every output has already
been collected~\cite{code-eval-entrypoints}.

\begin{enumerate}[label=\arabic*.]
\item \textbf{Topology embeddings.}
\code{model.eval_topology_embeddings} collects sixteen post-step raw embeddings
with all support switches on and Digital Corpora synthetic test traffic for
each of V0--V4. It stores one JSON per topology. When all five compatible
collections are present, it jointly projects eighty vectors with t-SNE.
These are temporal observations of five selected topologies, not eighty
separately generated topologies. The model hash, encoder, seed, and sample
count must match across collections. A plot-only path redraws saved records.

\item \textbf{Traffic embeddings.}
\code{model.eval_traffic_embeddings} uses V0 and collects the first post-step
all-on embedding with nonzero delivered traffic from each of T0, T1, Megalodon
T2, T3, T4, and TrEnvs~2--5. The resulting nine raw vectors are jointly
projected. Both embedding workflows use seeded two-dimensional t-SNE with
random initialization and perplexity
$\min(30,\max(1,(n-1)/3))$. Raw vectors and projection coordinates are saved.
Source normalization is not applied to the vectors passed to the t-SNE plot.

\item \textbf{Dataset comparison.}
\code{model.eval_datasets} evaluates the three frozen models and matched
baseline on V0. The source group contains \code{dc_t0}, \code{dc_pcap},
\code{mg_t0}, and \code{mg_pcap}; the load group contains T1, T3, and T4;
the congestion group contains TrEnvs~2--5. Three metrics per group give nine
bar figures. Reports retain raw repetition measurements and model/dataset
hashes.

\item \textbf{Topology comparison.}
\code{model.eval_topologies} evaluates the currently started member of V0--V4
using Digital Corpora test PCAPs and the network's host-count dataset.
After five compatible collections it produces three charts, each with five
topology groups and three model bars. For each topology it also writes three
CSV/PNG reward curves, using the mean raw return for each episode index across
repetitions. Five topologies therefore yield fifteen CSVs and fifteen reward
plots. These curves are evaluation returns and do not reproduce the paper's
combined, min--max-normalized training/testing Figure~9. Aggregate plot-only
mode checks the same model hashes, settings, seeds, and repetition budgets.

\item \textbf{Pruning comparison.}
\code{model.eval_pruning} reads the already generated training CSVs for
\code{dc_t0}, \code{dc_t0_no_pruning}, and \code{dc_t0_static_pruning}.
It plots raw episode rewards and a relative difference from the no-pruning
curve. Episode indices are aligned, and the relative statistic is
\begin{equation}
 100\frac{C_{\mathrm{variant}}-C_{\mathrm{none}}}
             {|C_{\mathrm{none}}|}.
\end{equation}
A zero no-pruning reward leaves this statistic undefined. This entry point
does not itself train the three variants or measure network communication
bandwidth.

\item \textbf{Enabled controllable switches.}
\code{model.eval_switches} runs one episode per model and the all-on baseline
with a specified network, test dataset, seed, and exact \code{--steps} horizon.
The point after decision $t$ is $\sum_i a_t^i$, and final counts appear in
the legend and JSON. These count support switches, not all switches. With
one distinct support per domain and no failures, total available switches
would be $|V|-N+\sum_i a_t^i$. Collection uses telemetry without delivery
captures. An explicitly longer horizon can continue into idle time after
the dataset finishes. The plot begins at the first decision rather than the
all-on reset state.
\end{enumerate}

\subsection{Verification performed for this source-based description}
\label{sec:verification}

The supplied paper was read through text extraction and visual inspection of
the learning equations and experimental-setup table. The methodology was
cross-checked against the preparation modules, scenario files, controller
handlers, numerical update code, evaluation modules, settings, and existing
unit fixtures. No full dataset generation, Mininet or Ryu launch, training
run, or evaluation run was performed in preparing this document.

The existing unit suites were executed in their designated runtimes:
24 model tests, 10 data-preparation tests, 9 network-preparation tests,
7 Mininet-module tests, and 6 controller tests, for 56 passing tests. They
exercise small fixtures and mocked clients, including embedding dimensions
and self-loop conventions, reward arithmetic, pruning-copy preservation,
checkpoint round trips, curriculum budgets, paired metrics, time stretching,
connectivity-safe switch selection, controller ownership, counter handling,
and OpenFlow message contracts~\cite{code-tests}.

These checks substantiate local contracts. They do not establish live replay
timing, end-to-end integration, learning convergence, physical power savings,
or equality with the original publication's numerical results. Those claims
require the separately executed integration and experimental protocols.

\clearpage
\section{Missing Components and Prospects for Faithful Replication}
\label{sec:missing}

\subsection{Assessment criteria}

This section concerns features and experiments described as part of EAGLE in
the paper, rather than every technique surveyed in its related work. The
collaborative average-reward objective, Actor--Critic networks, global reward
estimator, fixed-size graph embedding, curriculum stages, and magnitude-pruned
consensus are present. Likewise, the absence of an online MILP solver is not
itself a replication gap: the paper proposes RL as the practical alternative
to solving its formal optimization problem.

Three levels of reconstruction must be distinguished. A mathematical component
is \emph{algorithmically reconstructible} when its required computation is
sufficiently specified. A \emph{paper-consistent approximation} is possible
when the intended behavior is known but implementation or experimental choices
remain unspecified. \emph{Exact experimental replication is not recoverable}
when original configurations, datasets, deployment conditions, or records are
missing. These levels can apply simultaneously to different parts of one
feature. Table~\ref{tab:gaps} gives an overview, followed by the evidence and
qualification for each gap.

\begingroup
\small
\setlength{\LTpre}{0.5em}
\setlength{\LTpost}{0.5em}
\begin{longtable}{@{}P{0.25\textwidth}P{0.32\textwidth}P{0.36\textwidth}@{}}
\caption{Paper-to-code gaps and reconstruction feasibility. A missing completed
experiment is distinguished from a missing algorithm.}\label{tab:gaps}\\
\toprule
Paper component & Local status & Feasibility from the supplied material \\
\midrule
\endfirsthead
\toprule
Paper component & Local status & Feasibility from the supplied material \\
\midrule
\endhead
\midrule
\multicolumn{3}{r}{Continued on next page}\\
\endfoot
\bottomrule
\endlastfoot
Decentralized, time-varying AC-MARL network & Fixed local learner pool;
coordinator-mediated exchange; Ryu failover is separate & Static consensus
equations are implementable; original dynamic membership and distributed
recovery protocol are insufficiently specified. \\
Switch/interface shutdown and activation & Forwarding graph and rules change;
OVS processes/interfaces stay up & Functional actuation can be added;
original timing and device behavior cannot be recovered from the paper. \\
Balanced final domains and effective placement & Capacity is not enforced
after merging; placement is metadata with local endpoints & Constraints and
hop objective are known; final heuristic details and original deployment
are not uniquely recoverable. \\
Original topology classes & V1 does not preserve degree sequence; V4 count
differs; only five-topology scatter workflow & Required graph properties can
be recreated; exact graph ensemble and original realizations are unknown. \\
Original trace and TrEnv construction & Seeded window pooling, host-count
stretching, and fixed multipliers are local choices & Local inputs can be
regenerated; exact published traffic splits and schedules cannot be inferred. \\
Five competing frameworks & DeepConf, SmartFCT, Mystique, CollaQ, and CATE
implementations are absent & High-level descriptions are insufficient for
faithful implementations and trained benchmark policies. \\
Carbon-emissions comparison & Equation~(35), regional intensities, and CATE
comparison are absent & Arithmetic can be implemented; original regional
inputs and comparison results are unavailable locally. \\
Link/agent failure experiments & Reaction primitives and unit fixtures exist;
published experiment drivers and delay metrics do not & A new protocol is
possible; original schedules and result distributions are unspecified. \\
FABRIC / NSFNET evaluation & Deployment and experiment assets are absent &
Requires external testbed resources and missing original configurations. \\
Action-latency and communication-overhead study & Partial timing fields and
serialized-byte counts exist; full measurement suite is absent & Timers and
counting can be added; original hardware-specific results are not recoverable. \\
Embedding sensitivity / cross-validation & Configurable $g,e$ exist; sweep,
selection, and cross-validation workflow is absent & Embedding sizes are exact;
the original validation protocol and trained sweep results are unknown. \\
Complete published results and artifacts & Some local models/plots exist;
complete benchmark/result archive is absent & Local experiments can be run;
the publication's numerical claims cannot be established from partial outputs. \\
\end{longtable}
\endgroup

\subsection{Distributed AC-MARL and a time-varying communication network}
\label{sec:gap-distributed}

The paper's Section~V describes networked agents communicating with neighbors
through a time-varying network $cn_t$, exchanging local embeddings and actions
during operation and shared learning parameters during training. It also
describes removing failed agents while continuing collaborative operation.
The local implementation instead uses a fixed agent list and a central
coordinator that fetches every logical slot's telemetry, constructs the shared
state, and distributes complete batches. Its all-to-all consensus is a
uniform, static special case of the paper's weighting rule. Learner processes
do not communicate directly across controller hosts, and a failed learner
aborts its round~\cite{paper,code-workers,code-env}.

There is substantive \emph{physical controller} recovery: peer expiry,
ownership election, role switching, and HTTP endpoint fallback are present.
However, failed physical controllers retain their logical policy slots, whose
learning entities continue in the numerical runtime. This does not demonstrate
the paper's dynamic learning-agent membership or fully decentralized deployment.

The consensus equations are sufficient to implement weighted message mixing
over a supplied communication graph. They do not uniquely specify the
membership protocol, distributed state synchronization, failure-detector
interaction, or policy-slot treatment after agent removal. A distributed
extension consistent with the paper is feasible, but claiming a faithful
reconstruction of its original failure behavior would require additional
source or deployment evidence.

\subsection{Physical resource actuation and the scope of power savings}

The paper's control-loop description and overhead Figure~12 discuss
activation/deactivation of interfaces and OVS switches. The local action API
only changes support-switch availability in the active graph and reconciles
OpenFlow rules. There is no action-to-Mininet operation that shuts down a
bridge or its interfaces; the Mininet control socket exposes no such command
\cite{paper,code-routing,code-runtime}.

Logical removal provides a usable forwarding-level autoscaling experiment,
and Equation~\eqref{eq:power} assigns the associated modeled saving. It does
not establish that emulated or physical device consumption actually falls by
that amount. Implementing explicit bridge/interface transitions is possible
using the existing switch/port map and Mininet/OVS facilities, but consistent
reconnection, ownership, and telemetry handling would also need integration
work. The paper does not provide sufficient original device, command, or
hardware information to reproduce the published activation-latency distribution
exactly. This gap concerns actuation fidelity; the paper itself also uses a
simplified analytical power model.

\subsection{Final load balance and realized controller placement}

The balanced Louvain phase contains the intended capacity check and penalty,
but the final merge to $N$ communities does not enforce them. The inspected
V4 partition of 2, 94, and 53 switches exceeds its recorded move capacity of
50. The code penalizes a size difference \emph{greater than} two, whereas the
paper states a pairwise difference \emph{less than} two. Default move capacity
is topology-dependent, while the paper reports $\gamma=6$ for its training
network. Moreover, V0 preserves source domains rather than rerunning the
heuristic~\cite{paper,code-network,code-settings}.

The minimum-hop placement objective is implemented, but the default controller
processes communicate over loopback endpoints. Scenario placement switches do
not physically route management traffic through those switches or impose
corresponding controller-to-switch latency. Therefore the current emulator
does not by itself test the paper's claimed benefit from locating controllers
near their managed switches.

Hard final-capacity checks and a topology-aware control-plane deployment can
be constructed from the stated constraints and scenario graph. The paper
does not identify the replica's forced-merge extension or fully specify its
own tie-breaking, final community-count handling, and control-channel setup.
The intended objective is reproducible; its exact original realization is not
determined by these sources.

\subsection{Degree-preserving topologies and the original graph ensemble}

Class~1 in the paper's graph-embedding study and V1 in Table~III retain the
source degree distribution. The local random generator enforces only
connectivity, node count, and edge count. The inspected V1 has degrees up to
five, while the source's maximum degree is four. Thus the required
degree-preserving construction is absent, even though the V1 label and
$13$-node/$16$-link size are present~\cite{paper,code-network}.

A degree-preserving rewiring of the supplied V0 graph could reconstruct that
\emph{class property}. The exact published random graph cannot be recovered
without its randomization method, seed, or original edge list. V2 and V3
similarly match reported sizes without establishing identity of every original
edge or capacity. The source JSON records that V0 was reconstructed from an
earlier Mininet topology and Ryu port map; that provenance is useful but does
not independently prove equality with the paper's complete experimental setup.

GtsCe is available locally, but the supplied GraphML and V4 JSON represent
193 undirected links, compared with the paper's 386. Representing both
directions of those links would produce 386 arcs, so a counting convention is
a plausible explanation, not a verified conclusion. The sources do not justify
inventing 193 additional physical links. Finally, the implemented t-SNE
topology workflow samples sixteen times from each of five graphs; it does not
generate the larger collection of distinct graphs in the four classes shown
in the paper's Figure~4(a).

\subsection{Original traffic schedules and exact curriculum construction}

The code provides raw captures and reproducible local preparation rules, but
the paper does not specify the exact source-window split, window-to-host
assignment, host-count time stretch, trace boundary choices, or numerical
hotspot multipliers used in the original experiments. Its traffic embedding
discussion also describes extracting different trace intervals for unseen
load conditions, while this replica constructs T1, T3, and T4 by multiplying
the same test T0 windows~\cite{paper,code-data}.

The local 80/20 window split, synthetic payload conversion, $H$-dependent
stretching, and 1.5 regional multiplier are therefore explicit reconstruction
choices. TrEnv~5 and T3 coincide under the local rules, and intended hotspot
names are not evidence that particular original congestion levels have been
replicated. Training/test windows can also belong to the same original TCP
connection because splitting is not connection-wise.

It is possible to regenerate the \emph{local} manifests from the raw captures,
settings, and seeds. It is not possible to uniquely recover the publication's
train/test traffic, unseen trace intervals, or measured load distributions
without additional original manifests or records. Implementing another
paper-consistent traffic construction would be an approximation, not proof
of exact traffic fidelity.

\subsection{Published competing autoscaling frameworks}

The paper compares EAGLE with DeepConf, SmartFCT, Mystique, CollaQ, and CATE
and describes their different inputs, objectives, collaboration strategies,
and retraining requirements. None of these five algorithms is implemented
in the local source. The three default comparison models are EAGLE training
variants, and the all-on baseline is a no-autoscaling reference. These
comparisons cannot be described as reproducing the paper's state-of-the-art
benchmark comparisons~\cite{paper,code-evaluation,code-eval-entrypoints}.

The paper's overview is insufficient to recover all five complete algorithms,
their adapters to this SDN environment, hyperparameters, original reward
functions, and trained policies. Faithful implementations would require their
original papers or code and the original EAGLE comparison configurations.
A rough imitation from the overview would not meet the same standard. The
published training-time comparisons and benchmark-relative efficiency scores
are consequently not reconstructible from the supplied local implementation
and EAGLE paper alone.

\subsection{Carbon emissions and delivery metrics beyond completed-flow FCT}

The paper's Equation~(35) introduces regional carbon intensity and its
Figure~6 compares throughput, decisions, and carbon emissions with CATE.
The replica has a power model and throughput counters but no carbon-intensity
field, regional intensity dataset, carbon-emissions computation, CATE driver,
or corresponding figure generator~\cite{paper,code-reward,code-plots}.

The equation is sufficient to implement a carbon-cost calculation \emph{given}
the regional intensities and appropriate device/port coefficients. Those
original inputs and time-series records are not supplied. Assigning new
intensities could support a new experiment, but would not reproduce the
published emissions comparison.

Likewise, packet-delay measurement for the failure study is absent. The FCT
pipeline tracks successful delivery coverage, then drops undefined FCTs from
the recorded completed-flow events and aggregate means. There is no companion
reported loss rate, all-offered-flow completion ratio, or timeout-penalized FCT
metric in the comparison suite. Some such diagnostics could be added using
existing delivery records, but the original paper does not specify enough
packet matching and completion/censoring details to establish an identical
published FCT/delay estimator.

\subsection{Link and agent failure experiments}

Ryu reacts to port-state changes, rebuilds routes when the active graph
changes, and transfers physical ownership after a peer expires. Small unit
fixtures check disabled-support routing and controller ownership. The
published Figure~10 experiments are nevertheless absent as an executable
workflow: there is no prescribed failure-injection schedule, trial selector,
failed-link-count sweep, packet-delay collector, or agent-failure efficiency
comparison. The supervisor exposes restore/status, not an experiment-level
stop/failure command, and the Mininet socket has no link-failure operation
\cite{paper,code-runtime,code-routing,code-tests}.

A new experiment could use manual Mininet link changes, controlled process
termination, and the existing telemetry/reward measurements. The paper does
not provide exact failed-link choices, onset times, durations, recovery
schedules, or complete membership-handling details. Functional robustness can
therefore be investigated with a new, stated protocol, but its original
distributions and efficiency values are not recoverable from the supplied
material.

\subsection{FABRIC testbed and NSFNET generalization}

The paper reports a FABRIC deployment on NSFNET and a comparison with Mininet
and retrained alternatives in Figure~11. The local repository contains no
FABRIC slice-provisioning workflow, NSFNET scenario, distributed host setup,
testbed-specific traffic/measurement adapter, or original FABRIC result
archive. The controller launcher's ability to select an individual controller
and use reachable endpoints is a deployment primitive, not the reported
testbed experiment~\cite{paper,code-runtime}.

Recreating a FABRIC experiment requires testbed access, the missing topology
and resource mapping, link and traffic configurations, and a measurement
protocol. The paper also attributes some result variation to shared-link
background traffic, which cannot be reconstructed deterministically. A new
external deployment could test the same hypothesis; faithful reproduction of
the original deployment and measured background conditions is not possible
from these files alone.

\subsection{Overhead measurements and embedding sensitivity experiments}

The paper's Figure~12 separates inference, interface transitions, OVS
activation/deactivation, and route installation over 25 measurements.
The code records whole-step duration, rule-dispatch timestamps, a barrier
completion timestamp, serialized parameter sizes, and local IPC byte counts.
It does not provide the complete timing experiment, the OVS/interface
actuation path, an inference-only timing series, or a plotting workflow for
these latency distributions. Serialized contribution counts also do not
measure end-to-end distributed communication bandwidth~\cite{paper,code-routing,code-workers}.

Instrumentation could be extended to measure those components once explicit
resource actuation and distributed execution exist. The exact published
latencies additionally depend on hardware, system load, software versions,
control-channel placement, and timing boundaries that are not fully specified.
The original numerical overhead results cannot be derived from a static
source inspection.

The paper's Figure~13 and Section~VI-I discuss varying FeatherGraph $g$ and
$e$, cross-validation, and the selected representation's quality. The local
encoder supports those parameters and exactly determines embedding size
$4ge$ per domain. The repository lacks a dedicated parameter-sweep,
cross-validation, model-selection, and sensitivity-figure workflow. One can
reproduce the size relationship analytically and run a newly defined sweep;
the paper does not identify validation folds, budgets, or all settings needed
to reconstruct its selection experiment exactly.

The pruning mechanism and reward-curve comparison associated with Figure~14
\emph{are} implemented. Their presence should not be confused with a completed
distributed bandwidth study or with reproduction of the paper's original
learning curves. For example, the twenty-round ramp length and thirty policy
episodes are explicit local defaults rather than quantities recovered from
the publication.

\subsection{Original implementation choices and completeness of experimental evidence}

Several important details are defined by the local implementation but not
uniquely determined by the paper: exact FeatherGraph self-loop, feature,
frequency-grid, and boundary-node conventions; normalization calibration and
range; throughput/power reference values; Adam versus the schematic gradient
steps; batch and episode budgets; and the dynamic-pruning ramp length.
The current descriptor and losses can be reproduced faithfully from their
source. The paper alone does not establish that those choices yield the
same embeddings or policies as the original implementation. In particular,
the paper's switch-only graph description is less specific than the replica's
host-inclusive, boundary-overlapping observation graph
\cite{paper,code-embedding,code-learning,code-settings}.

The supplied paper also contains both general repeated estimator updates in
Section~V and a frozen-estimator setup in Section~VI-A. The replica's
pretrain/freeze procedure has direct support in the latter and is not marked
as a missing estimator implementation. Resolving every original update or
representation convention would require the original experiment's source
revision and configuration, rather than extrapolation from the equations.

Finally, completed artifacts and implemented workflows are different forms of
evidence. At inspection, the local output tree contained a named
\code{dc_t0} inference artifact, \code{dc_t0}/\code{dc_pcaps} training
checkpoints, traffic-embedding output, and a V0 topology-embedding collection.
It did not contain the complete three-model comparison archive, all five
topology collections, or the original benchmark/FABRIC/failure results.
This is an incomplete experimental record, not evidence that the existing
training and evaluation entry points are absent.

Additional local runs could populate the supported outputs using the stated
methodology. They would constitute experiments on this replica. Neither the
partial artifacts nor the passing unit fixtures substantiate the paper's
reported FCT improvements, power savings, zero-shot performance, robustness,
or training times. Recovering those original results requires missing
benchmark implementations and original experimental assets in addition to
executing the implemented core.

\clearpage
\begin{thebibliography}{99}
\small

\bibitem{paper}
D. Monaco, A. Sacco, F. Esposito, and G. Marchetto,
``Adaptive SDN Autoscaling via Generalizable Multi-Agent Reinforcement Learning
With EAGLE,'' \emph{IEEE Transactions on Network and Service Management},
vol.~23, pp.~6507--6522, 2026.
DOI: \href{https://doi.org/10.1109/TNSM.2026.3717361}{10.1109/TNSM.2026.3717361}.
Source: the PDF supplied with this repository.

\bibitem{code-snapshot}
EAGLE-replica local source tree, inspected 18 November 2026;
Git HEAD \code{37699c8af206e1a2daf3faae37e085efbaa9e5f7}.
Paths in the following entries are relative to the repository root.
Output-tracking policy: \path{.gitignore}.
Workflow context: \path{Documentation.md}, \path{README.md}, and
\path{pipeline.ipynb}. These guides had pre-existing local edits at inspection.

\bibitem{code-runtime}
Runtime and traffic execution:
\path{mininet_nets/launch_topology.py},
\path{mininet_nets/traffic.py},
\path{mininet_nets/traffic_flow.py},
\path{ryu_controller/launch_controllers.py},
\path{ryu_controller/requirements.txt},
\path{model/requirements.txt},
\path{mininet_nets/README.md}, and \path{ryu_controller/README.md}.

\bibitem{code-data}
Trace preparation and traffic variants:
\path{data_preparation/analyze_trace.py},
\path{data_preparation/load_and_split.py},
\path{data_preparation/prepare.py},
\path{data_preparation/replay.py},
\path{data_preparation/traffic_loads.py}, and
\path{data_preparation/traffic_congestions.py}.

\bibitem{code-network}
Topology construction, partitioning, and placement:
\path{network_preparation/prepare.py},
\path{network_preparation/balanced_louvain.py},
\path{network_preparation/placement.py},
\path{network_preparation/plot.py},
\path{network_preparation/source_v0.json}, and
\path{network_preparation/topology-zoo/GtsCe.graphml}.
Inspected generated examples: \path{networks/v0.json} through
\path{networks/v4.json}.

\bibitem{code-settings}
Defaults and validation: \path{config.py}, \path{settings/default.json},
\path{settings/no_pruning.json}, \path{settings/static_pruning.json}, and
\path{model/settings.py}.

\bibitem{code-routing}
Controller topology, routing, and telemetry:
\path{ryu_controller/controller.py},
\path{ryu_controller/topology_state.py},
\path{ryu_controller/routing.py},
\path{ryu_controller/openflow_rules.py}, and
\path{ryu_controller/telemetry.py}.

\bibitem{code-env}
Joint environment and runtime clients:
\path{model/joint_environment.py},
\path{model/controller_client.py},
\path{model/emulation.py}, \path{model/scenario.py}, and
\path{model/cli.py}.

\bibitem{code-datasets}
Dataset resolution, validation, and grouping: \path{model/datasets.py},
especially \code{load_dataset}, \code{Dataset.decisions}, and
\code{DATASET_GROUPS}.

\bibitem{code-reward}
Reward arithmetic and modeled watts: \path{model/reward.py},
\code{compute_reward}; coefficients in \path{settings/default.json}.

\bibitem{code-embedding}
Graph descriptors, calibration bounds, and collection:
\path{model/observations.py}, especially \code{embed_subnet},
\code{ObservationEncoder.normalize}, and \code{make_encoder};
\path{model/embeddings.py}.

\bibitem{code-learning}
MLPs, action selection, and gradient updates:
\path{model/actor_critic.py} and \path{model/learner.py},
especially \code{AgentLearner.learn}.

\bibitem{code-workers}
Learner process lifecycle and synchronized rounds: \path{model/workers.py},
especially \code{AgentPool.dispatch} and \code{AgentPool.run_round}.

\bibitem{code-consensus}
Pruning and parameter mixing: \path{model/consensus.py},
\code{pruning_fraction}, \code{pack_contribution}, and
\code{mix_contributions}.

\bibitem{code-training}
Training orchestration and curriculum budgets: \path{model/train.py} and
\path{model/settings.py}, especially \code{training_stages}.

\bibitem{code-checkpoints}
Persistence and provenance: \path{model/checkpoints.py},
\path{model/artifacts.py}, and \path{model/run_log.py}.

\bibitem{code-trials}
Episode aggregation and delivery analysis: \path{model/trials.py} and
\path{model/flow_records.py}, especially \code{run_trial},
\code{_managed_tcp}, and \code{_pcap_flows}.

\bibitem{code-evaluation}
Common frozen-policy comparison runner: \path{model/evaluate.py},
especially \code{evaluate_condition} and \code{make_environment}.

\bibitem{code-plots}
Metrics, bootstrap intervals, t-SNE, and plotting: \path{model/plots.py},
especially \code{compare_condition}, \code{interval}, and
\code{plot_embeddings}.

\bibitem{code-eval-entrypoints}
Evaluation entry points:
\path{model/eval_topology_embeddings.py},
\path{model/eval_traffic_embeddings.py},
\path{model/eval_datasets.py},
\path{model/eval_topologies.py},
\path{model/eval_pruning.py}, and
\path{model/eval_switches.py}.

\bibitem{code-tests}
Fixture and mocked-runtime test suites:
\path{model/tests/}, \path{data_preparation/tests/},
\path{network_preparation/tests/}, \path{mininet_nets/tests/}, and
\path{ryu_controller/tests/}. Test counts and outcomes reported here are
from execution during this inspection, not the publication's experiments.

\end{thebibliography}
\end{document}
